\begin{theorem}[5.10.5]
\EGAReaderTarget{ega4.c04.thm.5-10-5}{IV.2.s5.10.5-fr}
Sean $X$ un preesquema localmente noetheriano, $Z$ una parte de $X$ estable
por especialización y $\mathcal F$ un $\mathcal O_X$-módulo coherente. Para que
$\mathcal F$ sea $Z$-cerrado, es necesario y suficiente que se tenga
$\operatorname{prof}_Z(\mathcal F)\geq2$.
\end{theorem}

En virtud de (5.10.2), es posible limitarse al caso en que $\mathcal F$ es
$Z$-puro y $\operatorname{prof}_Z(\mathcal F)\geq1$. Además, decir que
$\operatorname{prof}_Z(\mathcal F)\geq2$ equivale a decir que, para toda parte
cerrada $Z_\alpha$ de $Z$, $\operatorname{prof}_{Z_\alpha}(\mathcal F)\geq2$;
y, del mismo modo, de (5.9.8) resulta que decir que $\mathcal F$ es $Z$-cerrado
equivale a decir que $\mathcal F$ es $Z_\alpha$-cerrado para todo $\alpha$.
Es posible, por tanto, limitarse ya al caso en que $Z$ es \emph{cerrado}. Puesto que la
cuestión es local, basta, para todo $x\in Z$, demostrar el teorema para
$\mathcal F|U$, siendo $U$ una vecindad abierta afín de $x$, y es posible por tanto
limitarse al caso en que $X=U$ es afín. Se sabe entonces que
$\operatorname{Ass}(\mathcal F)$ es finito (3.1.6), y como
$\operatorname{Ass}(\mathcal F)\subset X\setmin Z$, existe una sección $f$ de
$\mathcal O_X$ sobre $X$ tal que
$\operatorname{Ass}(\mathcal F)\subset X_f\subset X\setmin Z$ (\textbf{II}, 4.5.4); de ello se
deduce que $f$ es $\mathcal F$-regular (3.1.9) y que, para todo $y\in Z$, se
tiene $f_y\in\mathfrak m_y$, por lo que
$\operatorname{prof}(\mathcal F_y)
=1+\operatorname{prof}(\mathcal F_y/f_y\mathcal F_y)$
(\textbf{0}, 16.4.6).
La condición $\operatorname{prof}_Z(\mathcal F)\geq2$ equivale, por tanto, a
$\operatorname{prof}_Z(\mathcal F/f\mathcal F)\geq1$, o también, por (5.10.2),
al hecho de que $\mathcal F/f\mathcal F$ es $Z$-\emph{puro}, y basta ver que esta
última propiedad equivale al hecho de que $\mathcal F$ es $Z$-\emph{cerrado}.

Consideremos la sucesión exacta
$0\to\mathcal F\xrightarrow{f}\mathcal F
\to\mathcal F/f\mathcal F\to0$
(siendo inyectiva por hipótesis la homotecia de razón
$f:\mathcal F\xrightarrow{f}\mathcal F$); si se pone $W=X\setmin Z$, se tiene el diagrama
conmutativo
\[
\xymatrix{
0\ar[r] & \Gamma(X,\mathcal F)\ar[r]^f\ar[d] &
\Gamma(X,\mathcal F)\ar[r]\ar[d] &
\Gamma(X,\mathcal F/f\mathcal F)\ar[r]\ar[d] & 0\\
0\ar[r] & \Gamma(W,\mathcal F)\ar[r]_f &
\Gamma(W,\mathcal F)\ar[r] & \Gamma(W,\mathcal F/f\mathcal F)
}
\]
donde las dos filas son exactas (por ser $X$ afín). Si el homomorfismo de
restricción $\Gamma(X,\mathcal F)\to\Gamma(W,\mathcal F)$ es biyectivo, se
deduce de este diagrama que
\[
\Gamma(X,\mathcal F/f\mathcal F)
\longrightarrow\Gamma(W,\mathcal F/f\mathcal F)
\]
es inyectivo, y esto demuestra, por (5.9.8), que, si $\mathcal F$ es $Z$-cerrado,
$\mathcal F/f\mathcal F$ es $Z$-puro. Recíprocamente, supongamos que
$\mathcal F/f\mathcal F$ es $Z$-puro, y sea $s$ una sección de $\mathcal F$
sobre $W$; como $X_f\subset W$, existe un entero $n>0$ tal que
$f^n(s|X_f)$ se prolonga a una sección $t$ de $\mathcal F$ sobre $X$
(\textbf I, 1.4.1); además, como las restricciones de $t$ y de $f^ns$ a $X_f$ son
las mismas, de ello resulta que la restricción de $t$ a $W$ es igual a $f^ns$
en virtud de la relación $\operatorname{Ass}(\mathcal F)\subset X_f$
(5.10.2); como $f$ es $\mathcal F$-regular, bastará ver que $t$ es de la forma
$f^nt'$, donde $t'\in\Gamma(X,\mathcal F)$

\oldpage[IV-2]{116}

para demostrar que el homomorfismo
$\Gamma(X,\mathcal F)\to\Gamma(W,\mathcal F)$ es sobreyectivo y, por tanto,
biyectivo. Ahora bien, decir que $t=f^nt'$ significa que la imagen de $t$ en
$\Gamma(X,\mathcal F/f^n\mathcal F)$ es nula. Pero como
$f^k\mathcal F/f^{k+1}\mathcal F$ es isomorfo a
$\mathcal F/f\mathcal F$ y, por tanto, es $Z$-puro por hipótesis, se deduce de
(5.10.3), por inducción sobre $n$, que $\mathcal F/f^n\mathcal F$ es
$Z$-puro. Pero,
por definición, la imagen de $t|W=f^ns$ en
$\Gamma(W,\mathcal F/f^n\mathcal F)$ es nula, de donde se sigue la conclusión.

\begin{corollary}[5.10.6]
\label{IV.2.s5.10.6-fr}
Sea $\mathcal F$ un $\mathcal O_X$-módulo coherente. Para que $\mathcal F$ sea
$Z$-cerrado en un punto $x$, es necesario y suficiente que
$\operatorname{prof}_{Z_x}(\widetilde{\mathcal F}_x)\geq2$.
\end{corollary}

Esto resulta de (5.9.6) y (5.10.5).

\begin{theorem}[5.10.7]
\label{IV.2.s5.10.7-fr}
\emph{(Hartshorne).} Sean $X$ un preesquema localmente noetheriano y $Y$ una
parte cerrada de $X$. Supongamos que, para todo $y\in Y$, se tiene
$\operatorname{prof}(\mathcal O_{X,y})\geq2$; entonces, para toda componente
conexa $C$ de $X$, $C\setmin(C\cap Y)$ es conexo.
\end{theorem}

Es posible limitarse al caso en que $X$ es conexo; de (5.10.5) resulta entonces
que el homomorfismo canónico
$\mathcal O_X\to i_*(\mathcal O_X|X\setmin Y)$ (donde $i:X\setmin Y\to X$ es la inyección
canónica) es biyectivo. Por consiguiente, el homomorfismo de restricción
$\Gamma(X,\mathcal O_X)\to\Gamma(X\setmin Y,\mathcal O_X)$ es también biyectivo.
Basta ahora aplicar el lema (\textbf{III}, 7.8.6.1).
